\section{Jacobson spaces}
\label{section-space-jacobson}

\begin{definition}
\label{definition-space-jacobson}
Let $X$ be a topological space.
Let $X_0$ be the set of closed points of $X$.
We say that $X$ is {\it Jacobson} if every
closed subset $Z \subset X$ is the closure
of $Z \cap X_0$.
\end{definition}

\noindent
Note that a topological space $X$ is Jacobson if and only if
every nonempty locally closed subset of $X$
has a point closed in $X$.

\medskip\noindent
Let $X$ be a Jacobson space and let $X_0$ be the set
of closed points of $X$ with the induced topology.
Clearly, the definition implies that the morphism
$X_0 \to X$ induces a bijection between the closed
subsets of $X_0$ and the closed subsets of $X$.
Thus many properties of $X$ are inherited by $X_0$.
For example, the Krull dimensions of $X$ and $X_0$
are the same.

\begin{lemma}
\label{lemma-jacobson-check-irreducible-closed}
Let $X$ be a topological space. Let $X_0$ be the set
of closed points of $X$.
Suppose that for every point $x\in X$
the intersection $X_0 \cap \overline{\{x\}}$ is dense in $\overline{\{x\}}$.
Then $X$ is Jacobson.
\end{lemma}

\begin{proof}
Let $Z$ be closed subset of $X$
and $U$ be and open subset of $X$
such that $U\cap Z$ is nonempty.
Then for $x\in U\cap Z$ we have that $\overline{\{x\}}\cap U$ is a nonempty
subset of $Z\cap U$,
and by hypothesis it contains a point closed in $X$ as required.
\end{proof}

\begin{lemma}
\label{lemma-non-jacobson-Noetherian-characterize}
Let $X$ be a Kolmogorov topological space with a basis of quasi-compact
open sets.
If $X$ is not Jacobson, then there exists a non-closed point
$x \in X$ such that $\{x\}$ is locally closed.
\end{lemma}

\begin{proof}
As $X$ is not Jacobson there exists a closed set $Z$ and an open set $U$
in $X$ such that $Z \cap U$ is nonempty and does not contain points closed
in $X$. As $X$ has a basis of quasi-compact open sets we may replace $U$
by an open quasi-compact neighborhood of a point in $Z\cap U$ and so we may
assume that $U$ is quasi-compact open. By
Lemma \ref{lemma-quasi-compact-closed-point}, there exists a point
$x \in Z \cap U$ closed in $Z \cap U$,
and so $\{x\}$ is locally closed but not closed in $X$.
\end{proof}

\begin{lemma}
\label{lemma-jacobson-local}
Let $X$ be a topological space.
Let $X = \bigcup U_i$ be an open covering.
Then $X$ is Jacobson if and only if each $U_i$ is Jacobson.
Moreover, in this case $X_0 = \bigcup U_{i, 0}$.
\end{lemma}

\begin{proof}
Let $X$ be a topological space.
Let $X_0$ be the set of closed points of $X$.
Let $U_{i, 0}$ be the set of closed points of
$U_i$. Then $X_0 \cap U_i \subset U_{i, 0}$
but equality may not hold in general.

\medskip\noindent
First, assume that each $U_i$ is Jacobson.
We claim that in this case $X_0 \cap U_i = U_{i, 0}$.
Namely, suppose that $x \in U_{i, 0}$, i.e., $x$ is closed in
$U_i$. Let $\overline{\{x\}}$ be the closure
in $X$. Consider $\overline{\{x\}} \cap U_j$.
If $x \not \in U_j$, then $\overline{\{x\}} \cap U_j = \emptyset$.
If $x \in U_j$, then $U_i \cap U_j \subset U_j$
is an open subset of $U_j$ containing $x$.
Let $T' = U_j \setminus U_i \cap U_j$ and
$T = \{x\} \amalg T'$. Since $x$ is closed in $U_i \cap U_j$
(as $x$ is closed in $U_i$), we see that $T$, $T'$
are closed subsets of $U_j$ and $T$ contains
$x$. As $U_j$ is Jacobson we see that the closed points of
$U_j$ are dense in $T$. Because $T = \{x\} \amalg T'$
this can only be the case if $x$ is closed in $U_j$.
Hence $\overline{\{x\}} \cap U_j = \{x\}$. We conclude
that $\overline{\{x\}} = \{ x \}$ as desired.

\medskip\noindent
Let $Z \subset X$ be a closed subset (still
assuming each $U_i$ is Jacobson).
Since now we know that $X_0 \cap Z  \cap U_i
= U_{i, 0} \cap Z$ are dense in $Z \cap U_i$
it follows immediately that $X_0 \cap Z$ is
dense in $Z$.

\medskip\noindent
Conversely, assume that $X$ is Jacobson.
Let $Z \subset U_i$ be closed. Then
$X_0 \cap \overline{Z}$ is dense in $\overline{Z}$.
Hence also $X_0 \cap Z$ is dense in $Z$, because
$\overline{Z} \setminus Z$ is closed. As $X_0 \cap U_i
\subset U_{i, 0}$ we see that
$U_{i, 0} \cap Z$ is dense in $Z$.
Thus $U_i$ is Jacobson as desired.
\end{proof}

\begin{lemma}
\label{lemma-jacobson-inherited}
Let $X$ be Jacobson. The following types of subsets $T \subset X$
are Jacobson:
\begin{enumerate}
\item Open subspaces.
\item Closed subspaces.
\item Locally closed subspaces.
\item Unions of locally closed subspaces.
\item Constructible sets.
\item Any subset $T \subset X$ which locally on $X$
is a union of locally closed subsets.
\end{enumerate}
In each of these cases closed points of $T$ are
closed in $X$.
\end{lemma}

\begin{proof}
Let $X_0$ be the set of closed points of $X$. For any subset
$T \subset X$ we let $(*)$ denote the property:
\begin{itemize}
\item[(*)] Every nonempty locally closed subset of $T$ has a point
closed in $X$.
\end{itemize}
Note that always $X_0 \cap T \subset T_0$. Hence property $(*)$
implies that $T$ is Jacobson. In addition it clearly implies
that every closed point of $T$ is closed in $X$.

\medskip\noindent
Suppose that $T=\bigcup_i T_i$ with $T_i$ locally closed in $X$.
Take $A\subset T$ a locally closed nonempty subset in $T$,
then there exists a $T_i$ such that $A\cap T_i$ is nonempty, it is
locally closed in $T_i$ and so in $X$.
As $X$ is Jacobson $A$ has a point closed in $X$.
\end{proof}

\begin{lemma}
\label{lemma-finite-jacobson}
A finite Jacobson space is discrete.
A Jacobson space with finitely many closed points is discrete.
\end{lemma}

\begin{proof}
If $X$ is finite then the set $X_0 \subset X$ of closed points is finite.
Assume $X_0$ is is finite and $X$ is Jacobson.
Then $\overline{X_0} = X$ by the Jacobson property.
Now $X_0 =\{x_1, \ldots, x_n\} = \bigcup_{i = 1, \ldots, n} \{x_i\}$
is a finite union of closed sets, hence closed, so
$X = \overline{X_0} = X_0$. Every point is closed, and by
finiteness, every point is open.
\end{proof}

\begin{lemma}
\label{lemma-jacobson-equivalent-locally-closed}
\begin{slogan}
For Jacobson spaces, closed points see everything about the topology.
\end{slogan}
Suppose $X$ is a Jacobson topological space.
Let $X_0$ be the set of closed points of $X$.
There is a bijective, inclusion preserving correspondence
$$
\{\text{finite unions loc. closed subsets of } X\}
\leftrightarrow
\{\text{finite unions loc. closed subsets of } X_0\}
$$
given by $E \mapsto E \cap X_0$. This correspondence preserves
the subsets of locally closed, of open and of closed subsets.
\end{lemma}

\begin{proof}
We just prove that the correspondence $E \mapsto E \cap X_0$ is injective.
Indeed if $E\neq E'$ then without loss of generality $E\setminus E'$ is
nonempty, and it is a finite union of locally closed sets (details omitted).
As $X$ is Jacobson, we see that
$(E \setminus E') \cap X_0 = E \cap X_0 \setminus E' \cap X_0$ is not empty.
\end{proof}

\begin{lemma}
\label{lemma-jacobson-equivalent-constructible}
Suppose $X$ is a Jacobson topological space.
Let $X_0$ be the set of closed points of $X$.
There is a bijective, inclusion preserving correspondence
$$
\{\text{constructible subsets of } X\}
\leftrightarrow
\{\text{constructible subsets of } X_0\}
$$
given by $E \mapsto E \cap X_0$. This correspondence preserves
the subset of retrocompact open subsets, as well as complements
of these.
\end{lemma}

\begin{proof}
From  Lemma \ref{lemma-jacobson-equivalent-locally-closed} above,
we just have to see that if $U$ is open in $X$ then $U\cap X_0$ is
retrocompact in $X_0$ if and only if $U$ is retrocompact in $X$.
This follows if we prove that for $U$ open in $X$ then $U\cap X_0$ is
quasi-compact if and only if $U$ is quasi-compact.
From  Lemma \ref{lemma-jacobson-inherited} it follows that we may replace
$X$ by $U$ and assume that $U = X$.
Finally notice that any collection of opens $\mathcal{U}$ of $X$ cover
$X$ if and only if they cover $X_0$, using the Jacobson property
of $X$ in the closed $X\setminus \bigcup \mathcal{U}$ to find a point
in $X_0$ if it were nonempty.
\end{proof}



















